\documentclass[10pt,a4paper]{amsart}
\usepackage{amsmath,amssymb,amsthm,graphicx}
% Retained preamble for the Stacks Project chapter snapshot in this directory.
% The chapter's own preamble.tex is a copy of the shared Stacks Project preamble; the block
% below is sliced MECHANICALLY, line for line, from it. Nothing was added, removed or
% redefined; the wrapper class, the xr-hyper/hyperref wiring and the externaldocument list
% are omitted because the assembler supplies the document class.
%
% lines taken from preamble.tex: 3, 5, 9, 11, 12, 156, 157, 158, ... (51 lines total)
%
% package coverage restored from the chapter's own preamble (2026-10-13): xy, xy-2cell

\usepackage[all]{xy}
\xyoption{2cell}
\UseAllTwocells

\newenvironment{reference}{\comment}{\endcomment}
\newenvironment{slogan}{\comment}{\endcomment}
\newenvironment{history}{\comment}{\endcomment}
\theoremstyle{plain}
\newtheorem{theorem}[subsection]{Theorem}
\newtheorem{proposition}[subsection]{Proposition}
\newtheorem{lemma}[subsection]{Lemma}
\theoremstyle{definition}
\newtheorem{definition}[subsection]{Definition}
\newtheorem{example}[subsection]{Example}
\newtheorem{exercise}[subsection]{Exercise}
\newtheorem{situation}[subsection]{Situation}
\theoremstyle{remark}
\newtheorem{remark}[subsection]{Remark}
\newtheorem{remarks}[subsection]{Remarks}
\numberwithin{equation}{subsection}
\def\lim{\mathop{\mathrm{lim}}\nolimits}
\def\colim{\mathop{\mathrm{colim}}\nolimits}
\def\Spec{\mathop{\mathrm{Spec}}}
\def\Hom{\mathop{\mathrm{Hom}}\nolimits}
\def\Ext{\mathop{\mathrm{Ext}}\nolimits}
\def\SheafHom{\mathop{\mathcal{H}\!\mathit{om}}\nolimits}
\def\SheafExt{\mathop{\mathcal{E}\!\mathit{xt}}\nolimits}
\def\Sch{\mathit{Sch}}
\def\Mor{\mathop{\mathrm{Mor}}\nolimits}
\def\Ob{\mathop{\mathrm{Ob}}\nolimits}
\def\Sh{\mathop{\mathit{Sh}}\nolimits}
\def\NL{\mathop{N\!L}\nolimits}
\def\CH{\mathop{\mathrm{CH}}\nolimits}
\def\proetale{{pro\text{-}\acute{e}tale}}
\def\etale{{\acute{e}tale}}
\def\QCoh{\mathit{QCoh}}
\def\Ker{\mathop{\mathrm{Ker}}}
\def\Im{\mathop{\mathrm{Im}}}
\def\Coker{\mathop{\mathrm{Coker}}}
\def\Coim{\mathop{\mathrm{Coim}}}
\DeclareMathSymbol{\boxtimes}{\mathbin}{AMSa}{"02}
\def\QCohstack{\mathcal{QC}\!\mathit{oh}}
\def\Cohstack{\mathcal{C}\!\mathit{oh}}
\def\Spacesstack{\mathcal{S}\!\mathit{paces}}
\def\Quotfunctor{\mathrm{Quot}}
\def\Hilbfunctor{\mathrm{Hilb}}
\def\Curvesstack{\mathcal{C}\!\mathit{urves}}
\def\Polarizedstack{\mathcal{P}\!\mathit{olarized}}
\def\Complexesstack{\mathcal{C}\!\mathit{omplexes}}
\def\Pic{\mathop{\mathrm{Pic}}\nolimits}
\def\Picardstack{\mathcal{P}\!\mathit{ic}}
\def\Picardfunctor{\mathrm{Pic}}
\def\Deformationcategory{\mathcal{D}\!\mathit{ef}}
\usepackage{cleveref}
\crefname{originallemma}{original lemma}{original lemmas}
\Crefname{originallemma}{Original lemma}{Original lemmas}
\makeatletter
\newlabel{quot-lemma-pic-over-pic}{{0D26}{source}{}{sourcepointer-1}{}}
\newlabel{quot-lemma-pic-over-pic@cref}{{[originallemma][0][]0D26}{[1][0][]source}{}{}{}}
\makeatother
\crefname{originallemma}{original lemma}{original lemmas}
\Crefname{originallemma}{Original lemma}{Original lemmas}
\makeatletter
\newlabel{quot-situation-pic}{{0D25}{source}{}{sourcepointer-2}{}}
\newlabel{quot-situation-pic@cref}{{[originallemma][0][]0D25}{[1][0][]source}{}{}{}}
\makeatother
\crefname{originallemma}{original lemma}{original lemmas}
\Crefname{originallemma}{Original lemma}{Original lemmas}
\makeatletter
\newlabel{stacks-definition-gerbe-over-stack-in-groupoids}{{06P2}{source}{}{sourcepointer-3}{}}
\newlabel{stacks-definition-gerbe-over-stack-in-groupoids@cref}{{[originallemma][0][]06P2}{[1][0][]source}{}{}{}}
\makeatother
\crefname{originallemma}{original lemma}{original lemmas}
\Crefname{originallemma}{Original lemma}{Original lemmas}
\makeatletter
\newlabel{stacks-lemma-when-gerbe}{{06P1}{source}{}{sourcepointer-4}{}}
\newlabel{stacks-lemma-when-gerbe@cref}{{[originallemma][0][]06P1}{[1][0][]source}{}{}{}}
\makeatother
\crefname{originallemma}{original lemma}{original lemmas}
\Crefname{originallemma}{Original lemma}{Original lemmas}
\makeatletter
\newlabel{stacks-morphisms-definition-gerbe}{{06QC}{source}{}{sourcepointer-5}{}}
\newlabel{stacks-morphisms-definition-gerbe@cref}{{[originallemma][0][]06QC}{[1][0][]source}{}{}{}}
\makeatother
\title{The Picard stack is a gerbe over the Picard functor}
\author{The Stacks Project Authors}
\date{stacks-project:moduli.tex:local-snapshot-20261012; GNU FDL 1.2 or later}
\begin{document}
\maketitle
\noindent\textit{Editorial note.} Statement span is lines 1638-1642 of sources/stacks-moduli/pinned-original-source.tex; proof span is lines 1644-1673, the complete author proof (28 body lines).\par\medskip
\noindent\textit{Editorial note.} Navigation stubs: 5 cited-but-not-included labels are supplied as non-load-bearing pointers with their pinned tags where the snapshot tag index has them.\par\medskip
\noindent\textit{Editorial note.} extra\_packages is empty on purpose: the retained preamble already loads xy with options and a second declaration clashes in the assembler wrapper.\par\medskip
\noindent\textit{Editorial note.} Text is retained verbatim, including the author's own typos and placeholders.\par\medskip
\noindent\textit{Editorial note.} Queue index 11 of queue\_moduli.json (pb=1429, stmt 1638-1642, proof 1644-1673).\par\medskip
\noindent\textit{Editorial note.} stacks-moduli wave-1: fresh adjudication accept (2026-10-14); pool per Chief order 08:5x (dpa superseded to lane-s9); moduli.tex 53 labels/48 with proof judged 25 accept / 23 reject / 5 pb=0 env-dup rows excluded; spans cite-free and reference/slogan/history-free so no patchenv needed.\par\medskip
\noindent\textit{Editorial note.} Author writes 'This follows immediately from ...; here is a detailed explanation' - the detailed explanation (both conditions verified) is the proof and is fully in-span.\par\medskip
\noindent\textit{Allowed standard prerequisites (editorial).}\begin{itemize}
\item Explicit prerequisite boundary. The cited inputs are prerequisites; the argument in the retained span is not replaced by any cited result.
\end{itemize}\medskip
\ifdefined\hypertarget\hypertarget{sourcepointer-1}{}\fi
\noindent\textit{Source reference.} Original lemma 0D26: sources/stacks-moduli/pinned-original-source.tex, source lines 1673--1673. Cited by the retained statement or proof but not part of the collected spans; supplied only as a navigation pointer. The argument does not depend on its content.\par\medskip
\ifdefined\hypertarget\hypertarget{sourcepointer-2}{}\fi
\noindent\textit{Source reference.} Original lemma 0D25: sources/stacks-moduli/pinned-original-source.tex, source lines 1673--1673. Cited by the retained statement or proof but not part of the collected spans; supplied only as a navigation pointer. The argument does not depend on its content.\par\medskip
\ifdefined\hypertarget\hypertarget{sourcepointer-3}{}\fi
\noindent\textit{Source reference.} Original lemma 06P2: sources/stacks-moduli/pinned-original-source.tex, source lines 1673--1673. Cited by the retained statement or proof but not part of the collected spans; supplied only as a navigation pointer. The argument does not depend on its content.\par\medskip
\ifdefined\hypertarget\hypertarget{sourcepointer-4}{}\fi
\noindent\textit{Source reference.} Original lemma 06P1: sources/stacks-moduli/pinned-original-source.tex, source lines 1673--1673. Cited by the retained statement or proof but not part of the collected spans; supplied only as a navigation pointer. The argument does not depend on its content.\par\medskip
\ifdefined\hypertarget\hypertarget{sourcepointer-5}{}\fi
\noindent\textit{Source reference.} Original lemma 06QC: sources/stacks-moduli/pinned-original-source.tex, source lines 1673--1673. Cited by the retained statement or proof but not part of the collected spans; supplied only as a navigation pointer. The argument does not depend on its content.\par\medskip
\begin{lemma}
\label{lemma-pic-gerbe-over-pic-functor}
The morphism $\Picardstack_{X/B} \to \Picardfunctor_{X/B}$
turns the Picard stack into a gerbe over the Picard functor.
\end{lemma}

\begin{proof}
The definition of $\Picardstack_{X/B} \to \Picardfunctor_{X/B}$ being
a gerbe is given in Morphisms of Stacks, Definition
\ref{stacks-morphisms-definition-gerbe}, which in turn refers to
Stacks, Definition \ref{stacks-definition-gerbe-over-stack-in-groupoids}.
To prove it, we will check conditions (2)(a) and (2)(b) of
Stacks, Lemma \ref{stacks-lemma-when-gerbe}. This follows immediately from
Quot, Lemma \ref{quot-lemma-pic-over-pic}; here is a detailed explanation.

\medskip\noindent
Condition (2)(a).
Suppose that $\xi \in \Picardfunctor_{X/B}(U)$ for some scheme $U$ over $B$.
Since $\Picardfunctor_{X/B}$ is the fppf sheafification of the rule
$T \mapsto \Pic(X_T)$ on schemes over $B$
(Quot, Situation \ref{quot-situation-pic}), we see that there exists an
fppf covering $\{U_i \to U\}$ such that $\xi|_{U_i}$ corresponds
to some invertible module $\mathcal{L}_i$ on $X_{U_i}$.
Then $(U_i \to B, \mathcal{L}_i)$ is an object of
$\Picardstack_{X/B}$ over $U_i$ mapping to $\xi|_{U_i}$.

\medskip\noindent
Condition (2)(b). Suppose that $U$ is a scheme over $B$ and
$\mathcal{L}, \mathcal{N}$ are invertible modules on $X_U$
which map to the same element of $\Picardfunctor_{X/B}(U)$.
Then there exists an fppf covering $\{U_i \to U\}$
such that $\mathcal{L}|_{X_{U_i}}$ is isomorphic to $\mathcal{N}|_{X_{U_i}}$.
Thus we find isomorphisms between
$(U \to B, \mathcal{L})|_{U_i} \to (U \to B, \mathcal{N})|_{U_i}$
as desired.
\end{proof}

\end{document}
